% Some commands used in this file
\newcommand{\package}{\emph}

% \chapter{Introduction}

% \section{Sample Section}
% \label{sec:section}

% \subsection{Sample Subsection}

% \theoremstyle{plain}
% \theoremsymbol{}
% \newtheorem{Rule}[theorem]{Rule}


% \begin{Rule}
% \end{Rule}

% \begin{quote}
% \end{quote}

% \emph{}

% \begin{Rule}
%   Learn when to use \lstinline[showspaces=true]-\ - and
%   \lstinline-\@-.
% \end{Rule}

% Unless you use `french spacing', the space at the end of a sentence is
% slightly larger than the normal interword space.

% The rule used by \TeX{} is that any space following a period,
% exclamation mark or question mark is sentence-ending, except for
% periods preceded by an upper-case letter.  Inserting \lstinline-\-
% before a space turns it into an interword space, and inserting
% \lstinline-\@- before a period makes it sentence-ending.  This means
% you should write
% \begin{lstlisting}
% Prof.\ Dr.\ A. Steger is a member of CADMO\@.
% If you want to write a thesis with her, you
% should use this template.
% \end{lstlisting}

\chapter{Introduction}

One of the most relevant topics being currently being explored in the field of computer graphics is the efficient generation or recreation of high-fidelity 3D scenes based on their 2D representations in the form of photographs taken from different angles. Example applications include: 
\begin{enumerate}
    \item Improving the user experience of online marketplaces by allowing buyers to inspect a model of their item of interest based on the seller's uploaded pictures,
    \item Digital memorialization of cultural and historical heritage sites, in particular for regions affected by war or natural disasters,
    \item Creating interactive video game or simulation environments based on real world observations.
\end{enumerate}

Techniques that have demonstrated significant success in this area are primarily categorized under the concept of radiance fields. A radiance field can be mathematically defined as a function $ L: \mathbb{R}^3 \times S^2 \rightarrow \mathbb{R}^3 $, where $ \mathbb{R}^3 $ on the left side represents the three-dimensional world space of the scene, $ S^2 $ denotes the unit sphere of directions surrounding each point, and $ \mathbb{R}^3 $ on the right side corresponds to the color space. Formally, this can be expressed as:

$$L(x, y, z, \theta, \phi) = (R, G, B)$$

In this equation, $ (x, y, z) $ denotes the spatial coordinates of a point in the scene, while $ (\theta, \phi) $ represents the viewing direction, defined in spherical coordinates (polar and azimuthal angle respectively). 

The first major advancements in radiance fields, Neural Radiance Fields (NeRFs~\cite{NeRF}), represented them implicitly using a neural network. However not so long after did explicit solutions emerge as well, offering superior real-time capabilities~\cite{3DGS}.

%This framework allows for the representation of the light emitted from any point in the scene in a specific direction, facilitating the reconstruction of 3D structures from their 2D projections. 

There are various rendering techniques used by the above solutions. The first group, represented by rasterization only roughly approximates light transfer, but achieves superb results in terms of time. The second, based on ray-tracing, aims to model this complex phenomenon in a more physically-accurate way, which for a long time in graphics was not possible due to limitations in hardware, but is now becoming more widely adopted and explored.

In this thesis, we focus on a framework devised by Condor et al.~\cite{DSYG} and presented at the ACM SIGGRAPH 2025 conference that combines an explicit representation scene representation with physically-based rendering. We aim to extend this framework in two major ways. First, in terms of speed, with a custom path-tracer specialized for the kernel mixture model scene representation. Second, with further research of its capabilities to render materials with complex surfaces such as textiles, and fluid dynamics, among others. 

% 3) we offer a unified solution for both that can be faster while still being physically based, unbiased, and capable of modelling scattering, relighting (light transport). 

